\begin{center}
\textbf{\large{Abstract}}
\end{center}
\addcontentsline{toc}{chapter}{Abstract}

\begin{small}

Database management systems (DBMS) are critical infrastructure components, and their vulnerabilities represent serious security threats. Automated vulnerability detection through fuzzing is a promising approach, but existing grammar-based fuzzers apply uniform random sampling over production rules, wasting mutation budget on low-value syntactic paths and failing to prioritize the structural patterns most likely to expose bugs.

This thesis presents DBMS-Nautilus, a grammar-based greybox fuzzing system targeting SQLite. A structural primitives grammar of 526 rules encodes SQL patterns known to trigger vulnerability classes --- including integer overflow in expression evaluation, use-after-free in query planner optimizations, and null pointer dereference in column resolution --- without hardcoding specific proof-of-concept inputs. The grammar was iteratively refined based on evidence from coverage analysis and crash investigation.

Experiments across 4 SQLite versions containing 6 known CVEs show that DBMS-Nautilus rediscovers 4 of the 5 detectable target CVEs through compositional grammar generation, and discovers 11 distinct bug classes --- including 3 previously unreported defects --- across 45 unique crash instances. A comparative evaluation against a generic EBNF baseline grammar demonstrates that the proposed grammar discovers $79\times$ more unique root causes per campaign despite $2.6\times$ lower throughput, indicating that grammar design is the dominant factor for vulnerability discovery.

\vspace*{1cm}
\textbf{Keywords}: Grammar-based Fuzzing, Greybox Fuzzing, DBMS Security, Vulnerability Detection, SQLite
\end{small}
